\chapter{Reproducibilidad}\label{ap:repro}

Este apéndice documenta el entorno de ejecución y los comandos exactos que
regeneran cada familia de resultados de la memoria. Todos los comandos se lanzan
desde la raíz del repositorio y operan sobre el panel estándar descrito en el
Capítulo de metodología.

\section{Entorno de ejecución}

El proyecto se ejecuta sobre \texttt{Python 3.11} (la dependencia
\texttt{pytorch-wavelets} no es compatible con Python 3.12). Las dependencias se
instalan desde \texttt{requirements.txt} en la raíz del repositorio mediante
\texttt{pip install -r requirements.txt}. Las librerías principales son:

\begin{itemize}
  \item \textbf{Cómputo numérico y datos}: \texttt{numpy}, \texttt{scipy},
        \texttt{pandas}, \texttt{pyarrow}.
  \item \textbf{Modelos clásicos (complexity ladder)}: \texttt{scikit-learn}
        (Ridge, Lasso, ElasticNet) y \texttt{lightgbm}.
  \item \textbf{Modelo profundo (Stockformer)}: \texttt{torch},
        \texttt{PyWavelets}, \texttt{pytorch-wavelets}, \texttt{tensorboard}.
  \item \textbf{Grafo dinámico}: \texttt{networkx}, \texttt{gensim},
        \texttt{fastdtw} y \texttt{GraphEmbedding} (Struc2Vec).
  \item \textbf{Estadística e ingeniería de cartera}: \texttt{statsmodels} y
        \texttt{cvxpy} para los pesos cost-aware.
  \item \textbf{Datos de mercado}: \texttt{yfinance} y \texttt{lxml}
        (lista de tickers del S\&P 500 vía \texttt{pandas.read\_html}).
  \item \textbf{Aplicación e informes}: \texttt{matplotlib},
        \texttt{streamlit}, \texttt{plotly}, \texttt{tqdm}, \texttt{psutil}.
\end{itemize}

Para entrenamiento en GPU (Colab o Linux) se instala \texttt{torch} con la rueda
CUDA correspondiente, según se indica en la cabecera de \texttt{requirements.txt}.
El resto del pipeline (ladder de CPU, estrategia semanal y robustez) corre
íntegramente en CPU.

\section{Semillas, determinismo y split estándar}

Toda la evaluación se ancla a un \emph{split temporal fijo} que garantiza que
todos los modelos comparados vean exactamente las mismas entradas; solo cambia
el modelo. El panel estándar se construye en \texttt{lib/data\_panel.py} con:

\begin{itemize}
  \item \textbf{Offset de etiquetas}: 60 filas (\texttt{ALPHA360\_LAG = 60}),
        el buffer del lag de Alpha360. Es el único invariante anti-fuga,
        verificado por \texttt{scripts/audit\_alignment.py}.
  \item \textbf{Variables}: el conjunto base \emph{F-base}
        ($F\!\approx\!375$ regresores: Alpha360 más indicadores técnicos),
        excluyendo las variables \texttt{MACRO\_*} para el ranking.
  \item \textbf{Split temporal fijo}: 0.75 / 0.125 / 0.125
        (entrenamiento / validación / test), con normalización ajustada
        únicamente sobre el tramo de entrenamiento.
\end{itemize}

La semilla maestra es \texttt{SEED = 0} y el número de réplicas bootstrap para
los intervalos de confianza del IC es \texttt{N\_BOOT = 2000}
(ver \texttt{scripts/run\_cpu\_ladder.py} y \texttt{scripts/run\_ladder\_analysis.py}).
El long-short se evalúa sobre el decil extremo (\texttt{QUANTILE = 0.1}) con una
comisión de 10 puntos básicos por lado (\texttt{FEE = 0.001}). Con estas
constantes fijadas, los resultados son reproducibles ejecución tras ejecución.

\section{Comandos de reproducción}

\subsection{Complexity ladder}

Evalúa todos los modelos de CPU (sanity, lineales y árboles) sobre el panel
estándar y escribe \texttt{results/ladder\_results.csv} más el IC diario por
modelo en \texttt{results/ladder\_daily\_ic/}:

\begin{verbatim}
python scripts/run_cpu_ladder.py \
    --data_dir data/Stock_SP500_2018-01-01_2026-03-16
\end{verbatim}

El análisis integra el Stockformer entrenado (desde su salida de inferencia
guardada) como peldaño superior, ejecuta los tests de significancia pareados y
genera \texttt{results/ladder\_full\_results.csv},
\texttt{results/significance\_tests.csv} y las figuras en
\texttt{results/figures/}:

\begin{verbatim}
python scripts/run_ladder_analysis.py
\end{verbatim}

\subsection{Estrategia semanal market-neutral (Tier A)}

Pipeline cost-aware de extremo a extremo (panel semanal $\rightarrow$ ensemble
LightGBM + ElasticNet $\rightarrow$ neutralización beta $\rightarrow$ suavizado
EWMA $\rightarrow$ pesos cost-aware con \texttt{cvxpy} $\rightarrow$ vol-target
$\rightarrow$ overlay de régimen). Produce
\texttt{results/weekly\_strategy\_summary.csv} y la curva de equity:

\begin{verbatim}
python scripts/run_weekly_strategy.py
\end{verbatim}

Los principales parámetros de construcción son configurables:

\begin{verbatim}
python scripts/run_weekly_strategy.py \
    --cost_bps 8 --gross 2 --name_cap 0.04
\end{verbatim}

\subsection{Robustez walk-forward y atribución}

Reentrena el ensemble periódicamente, concatena $\sim$215 semanas
out-of-sample en una única curva y descompone el Sharpe neto por etapa de
construcción. Genera \texttt{results/weekly\_robustness\_summary.csv} y
\texttt{results/weekly\_attribution.csv}:

\begin{verbatim}
python scripts/run_weekly_robustness.py --init_train 200 --step 26
\end{verbatim}

Las variantes con datos adicionales (fundamentales Tier B y realized-vol Tier C)
escriben sus propios resúmenes
(\texttt{weekly\_robustness\_summary\_fund.csv} y
\texttt{weekly\_robustness\_summary\_realized.csv}):

\begin{verbatim}
python scripts/run_weekly_robustness.py --with_fundamentals
python scripts/run_weekly_robustness.py --with_realized
\end{verbatim}

\subsection{Generación de tablas LaTeX}

Una vez generados los CSV anteriores, el siguiente comando reconstruye todos los
fragmentos LaTeX de tablas en \texttt{MEMORIA/tfm/tablas/} a partir de
\texttt{results/*.csv}:

\begin{verbatim}
python3 scripts/build_latex_tables.py
\end{verbatim}

Las tablas completas resultantes se recogen en el Apéndice~\ref{ap:tablas}.
